% Generated draft: requires mathematical and source review.
\documentclass[11pt]{article}
\usepackage[margin=25mm]{geometry}
\usepackage{fontspec}
\setmainfont{Latin Modern Roman}
\usepackage{amsmath,amssymb,mathtools,unicode-math}
\setmathfont{Latin Modern Math}
\usepackage{enumitem,needspace,xurl,hyperref}
\setlength{\parindent}{0pt}
\setlength{\parskip}{5pt}
\setlength{\emergencystretch}{3em}
\newcommand{\N}{\mathbb{N}}
\newcommand{\Z}{\mathbb{Z}}
\newcommand{\Q}{\mathbb{Q}}
\newcommand{\R}{\mathbb{R}}
\newcommand{\C}{\mathbb{C}}
\newcommand{\F}{\mathbb{F}}
\newcommand{\A}{\mathbb{A}}
\newcommand{\PP}{\mathbb{P}}
\newcommand{\E}{\mathbb{E}}
\newcommand{\eps}{\varepsilon}
\newcommand{\dd}{\,\mathrm{d}}
\begin{document}
\section*{UnsolvedMath ID: 30003453}
\textbf{Source problem code:} catalogue entry
\subsection{Definitions and mathematical statement}
Let \(n\geq 1\) denote the quantity called the sample size in the supplied statement, let \(\widehat f_n\) denote its maximum-likelihood boundary estimator, and write \(K_n\) for the number of observations described there as ``support points on the MLE.'' These are provisional notations: the observation model, admissible boundary functions, existence and selection of the maximum-likelihood estimator, and the precise criterion for counting support points are not specified in the record. The stated correction has magnitude
\[
\Delta_n=\frac{K_n}{n}.
\]
\textbf{Question as supplied.} In the irregular nonparametric support-boundary model, does a Bayesian procedure automatically produce the finite-sample correction \(\Delta_n\) to the maximum-likelihood estimator?

A precise mathematical assertion remains unresolved: the record does not identify the prior, the posterior summary or decision rule, the estimand to which the correction applies, or the sign and manner of applying the correction. In particular, it does not specify whether the target is the boundary itself or a scalar functional of it, whether the assertion concerns an exact finite-sample identity or an approximation, or over which models and Bayesian procedures it is quantified. Thus the displayed correction must not be read as an asserted identity between a Bayesian estimator and a corrected maximum-likelihood estimator.

The supplied dated literature assessment labels the problem partially solved. It reports automatic bias correction for compound-Poisson priors in specified classes, failures for linear boundaries, and no known universal automatic correction theorem. Those classes, failure conditions, and correction formulas are not given here; these are supplied curation claims rather than independently verified mathematical conclusions.
\subsection{Short English statement}
The question asks whether Bayesian inference automatically incorporates a correction of size \(K_n/n\), where \(K_n\) counts support points on a maximum-likelihood boundary estimate. The statistical model, prior, target, correction direction, and quantifiers require clarification before this can become a precise theorem or conjecture.
\subsection{Sources}
\begin{itemize}
\item Supplied local UnsolvedMath problem registry: permanent ID \(30003453\), exact source problem code \texttt{catalogue entry}, title ``Bayesian Bias Correction in Support-Boundary Models,'' in the Oberwolfach Reports Open Problems collection. The record gives proposed year 2017 and names no individual proposer.
\item Supplied source citation: \emph{Statistical Recovery of Discrete, Geometric and Invariant Structures} (2018), Oberwolfach Reports, pp.\ 949--999, DOI \texttt{10.4171/owr/2017/16}. Original Oberwolfach report: \url{https://doi.org/10.4171/owr/2017/16}. The supplied background describes these reports as documenting mathematical workshops and their open-problem discussions. The citation's 2018 date and its DOI's 2017 component are retained as supplied; their relationship has not been established here.
\item Supplied literature source: \emph{Nonparametric Bayesian analysis of the compound Poisson prior}, \url{https://arxiv.org/abs/1809.04140}. No authors or theorem references are supplied in the record.
\item Literature assessment checked at \texttt{2026-08-21}: supplied dated curation, not independently verified. Its status is ``partially solved''; it reports correction results for compound-Poisson priors in specified classes, failures for linear boundaries, and no known universal automatic correction theorem. The background explicitly calls this a dated literature triage derived from the supplied OWR curation and requests independent verification of time-sensitive status and cited claims.
\item The registry labels the statement \texttt{corrected\_verified} and records a manual semantic, typographic, and LaTeX review claiming self-containedness and preservation of the target. This is supplied review metadata, not independent verification of the formulation or mathematical correctness.
\end{itemize}
\end{document}
